% TOP_PROBLEM: {"id": 20002498, "title": "Non-self-adjoint spectral theory: completeness, reconstruction, and matrix-function bounds", "category_id": 16, "category": {"id": 16, "name": "physics", "display_name": "Mathematical Physics", "description": "Problems at the intersection of mathematics and physics.", "slug": "physics", "order_index": 16, "created_at": "2026-07-31T15:26:25.671Z"}, "status": "partially_solved"}
% FIRST_POSTED: null
% ENTRY_KIND: statement_only
% !TeX program = lualatex
\documentclass[11pt]{article}
\usepackage[margin=25mm]{geometry}
\usepackage{fontspec}
\setmainfont{Latin Modern Roman}
\usepackage{amsmath,amssymb,mathtools}
\usepackage{unicode-math}
\setmathfont{Latin Modern Math}
\usepackage{xurl,hyperref}
\hypersetup{colorlinks=true,urlcolor=blue}
\setcounter{secnumdepth}{0}
\setlength{\parindent}{0pt}
\setlength{\parskip}{5pt}
\setlength{\emergencystretch}{3em}
\newcommand{\R}{\mathbb R}
\newcommand{\N}{\mathbb N}
\newcommand{\Z}{\mathbb Z}
\newcommand{\Q}{\mathbb Q}
\newcommand{\C}{\mathbb C}
\newcommand{\D}{\mathbb D}
\newcommand{\F}{\mathbb F}
\newcommand{\sref}[2]{\hyperlink{source:#1:#2}{[S#2]}}
\newcommand{\cataloguescope}[1]{#1}
\begin{document}
% UnsolvedMath ID 20002498; source catalogue code AIM-PHYSICS-0006
\setcounter{section}{20002497}
\section{Non-self-adjoint spectral theory: completeness, reconstruction, and matrix-function bounds}\label{20002498}
{\small UnsolvedMath ID 20002498 · Mathematical Physics}
\subsection{Definitions and mathematical statement}

\paragraph{Shared notation.}$\mathbb N$, $\mathbb Z$, $\mathbb Q$, $\mathbb R$, and $\mathbb C$ denote the natural numbers, integers, rationals, reals, and complex numbers, respectively. All additional parameters and hypotheses are specified in the question.


\paragraph{Statement recorded in the supplied source.}
This collection asks the following questions about non-self-adjoint spectral theory.
\begin{enumerate}
\item For $A_\alpha=-d^2/dx^2+i|x|^\alpha$ on the real line, or on the positive half-line with a Dirichlet condition at zero, do the eigenfunctions form a complete system when $0<\alpha\leq2/3$?
\item On $\mathbb R^2_+=\{(x,y):y>0\}$, consider
\[A=-\partial_x^2-(\partial_y-ix^2/2)^2+icy,\qquad D(A)=H^1_0(\mathbb R^2_+)\cap\{u:Au\in L^2(\mathbb R^2_+)\}.\]
Is its spectrum nonempty for every parameter $c$?
\item Let $f$ be continuous, complex-valued, two-periodic, odd, and symmetric about $1/2$. Put $f_n(x)=f(nx)$, and define $A$ initially by $A(\sqrt2\sin(n\pi x))=f_n$. For $p>1$, $p\neq2$, find necessary and sufficient conditions on the Fourier coefficients of $f$ for $A$ to be boundedly invertible on $L^p(0,1)$, equivalently for $(f_n)$ to be a Schauder basis.
\item In $u'(t)=Au(t)$ with $A=L+B$, where $L$ is known self-adjoint and $B$ is unknown non-self-adjoint, recover $A$ or its discrete spectrum from observations $\omega(t)=\langle u(t),g\rangle$. How should the initial state and observation vector be chosen to detect all eigenvalues and generalized eigenvectors? What convergence and smoothness suffice for the spectral expansion? Can the complex eigenvalues and their multiplicities be recovered in finite time, with the smallest possible finite number of observations?
\item Under what hypotheses can the instability index of a non-self-adjoint differential operator be determined from a finite matrix obtained by restriction to trigonometric polynomials? The instability index counts unstable eigenvalues with multiplicity; the question includes controlling truncation errors and the required matrix size.
\item For a bounded operator on a Banach space, characterize when $\|A\|-\operatorname{spr}(A)$ is zero or strictly positive. Determine the best factor $m\in(0,1]$ in $\operatorname{spr}(A)\leq m\|A\|$ for the operator classes under consideration, including the relevance to discrete spectra of compact perturbations.
\item For a nonnormal $n\times n$ matrix, set
\[s(A,\varepsilon)=\inf_{V^{-1}(A+\Delta)V\ \mathrm{diagonal}}\bigl(\varepsilon\|V\|\|V^{-1}\|+\|\Delta\|\bigr).\]
Does a constant $C_n$ independent of $A$ give $s(A,\varepsilon)\leq C_n\sqrt\varepsilon$?
\item For a bounded operator $A$ with numerical range $W(A)$, is there a universal $C\geq2$ such that $\|f(A)\|\leq C\max_{z\in W(A)}|f(z)|$ for every analytic $f$ defined on a neighborhood of that range? Can $C\leq11.08$ be used?
\end{enumerate}
\par\smallskip\textit{Formulation references:} \sref{20002498}{1}, \sref{20002498}{2}.
\subsection{Short English statement}
Resolve the listed completeness, reconstruction, stability, approximation, and numerical-range questions for non-self-adjoint operators.
\subsection{Sources}
\begin{itemize}
\item[\hypertarget{source:20002498:1}{S1}] ULAM.AI and the UnsolvedMath Contributors, UnsolvedMath: A Curated Collection of Open Mathematics Problems, record 20002498 (AIM-PHYSICS-0006). Catalogue attribution under CC BY 4.0.
\url{https://huggingface.co/datasets/ulamai/UnsolvedMath}.
\item[\hypertarget{source:20002498:2}{S2}] Original problem formulation and mathematical context.
\url{https://aimath.org/pastworkshops/nonselfadjointproblems.pdf}.
\end{itemize}
\label{end:20002498}
\end{document}
